%! Function Model Construction and Application — Unit 1, Lesson 1.14 — Student Packet Cover Sheet
\documentclass[10pt]{article}
\usepackage{precalculus-article}
\usepackage{precalculus-boxes}
\usepackage{ltablex}
\keepXColumns

\begin{document}

% Full-bleed navy banner
\begin{tikzpicture}[remember picture, overlay]
  \fill[navy] (current page.north west)
    rectangle ([yshift=-0.9in]current page.north east);
\end{tikzpicture}
\vspace{-0.8in}

\begin{center}
  {\color{white}\textbf{\LARGE AP Precalculus}}\\[4pt]
  {\color{skymid}\large\textbf{Unit 1: Polynomial and Rational Functions}}\\[6pt]
  {\color{white}\textbf{\large Lesson 1.14 \quad Function Model Construction and Application}}
\end{center}

\vspace{0.22in}
\namedateperiod
\vspace{0.18in}

\begin{learningtargetbox}
\small
\begin{enumerate}[leftmargin=1.5em, itemsep=5pt]
  \item I can construct a polynomial or piecewise function model from a contextual description
        using restrictions, transformations, or regression techniques. \hfill\textit{(LO 1.14.A)}
  \item I can apply a function model to answer questions about predicted values, rates of change,
        and average rates of change, including appropriate units. \hfill\textit{(LO 1.14.C)}
\end{enumerate}
\end{learningtargetbox}

\vspace{0.14in}

\begin{tocbox}
\small
\renewcommand{\arraystretch}{1.5}
\begin{tabularx}{\linewidth}{@{} c l X r @{}}
\rowcolor{sky}
\textbf{\#} & \textbf{Component} & \textbf{Description} & \textbf{Score} \\
\midrule
1 & Warm-Up        & Spiral review: average rate of change, evaluating polynomials, identifying model type & \blank{1.2cm} \\
2 & Group Activity & Engineering Design Challenge: construct and apply models in three scenarios & \blank{1.2cm} \\
3 & Exit Ticket    & Evaluate a quadratic model, state a domain restriction, interpret rate of change & \blank{1.2cm} \\
4 & Homework       & Model construction, rational models, rates of change with units, piecewise pricing & \blank{1.2cm} \\
\midrule
  & \multicolumn{2}{r}{\textbf{Total}} & \blank{1.2cm} \\
\end{tabularx}
\end{tocbox}

\end{document}
